\documentclass[]{../../../../tex/classes/styledArticle}
\usepackage{../../../../tex/packages/hebrewSupport}
\usepackage{../../../../tex/packages/mathShortcuts}
\usepackage{../../../../tex/packages/theoremsSupport}

\author{שחר פרץ}
\title{חדו''א 2א $\sim$ \textit{הרצאה 10}}
\begin{document}
	\maketitle
	\textbf{מרצה: }יעל קרשון
	
	\section*{סטרלינג}
	\[ n! \sim \sqrt{2n\pi} \cl{\frac{n}{e}}^{n} \]
	כלומר: 
	\[ \frac{n!}{\sqrt{2 \pi n}\cl{\frac{n}{e}}^{n}} \rrr{n \to \infty} 1 \]
	
	מסתמך על השוואה לפונקציית גאמא (האינטגרל $\Gamma(n) = \int^{\infty}_0 t^{n}e^{-t}\dt$), שינויי משתנים וטור טיילור, ובסוף מז'וראנטה. הפרטים ברשימות של שירי. 
	
	להרבה שימושים מספיקה יותר קלה של סטרלינג, והיא ש־$\sqrt[n]{n!} \sim \frac{n}{e}$ כאשר $n$ שואף ל־$\infty$. 

%	TODO
	\section*{המשך רציפות במ''ש}
	\subsection*{החלפת גבול ונגזרת}
	\theo{נתונות $f, gf_n \co [a, b] \to \R$ ונניח ש־$f_n \in C^{1}([a, b])$ נקודתית וכן $f_n \to f, f_n' \to g$ במ''ש. אז $f_n \to f$ במ''ש וגם $f' = g$. }\begin{proof}
		$f_n \in C^{1}([a, b])$ ולכן $f_n'$ רציפות ובפרט אינטגרביליות. $g$ גבול במ''ש של פונקציות רציפות, ולכן $g$ רציפה אזי בעצמה אינטגרבילית. נקבע $x_0 \in [a, b]$ כלשהו. נגדיר: 
		\[ G(x) = \int^{x}_{x_0} \dequad \ g(t)\dt \quad\quad h_n(x) = \int^{x}_{x_0}\dequad \ f_n'(t)\dt \]
		ממשפט על גבול תחת האינטגרל, מקבלים $h_n(x) \to G(x)$ במ''ש (הוכחנו בהרצאה הקודמת). מניוטון לייבניץ קיבלנו: 
		\[ \underbrace{h_n(x)}_{\overset{u}{\to}G(x)} + \underbrace{f_n(x_0)}_{\overset{u}{\to}f(x_0)} = f_n(x) \]
		לכן $f_n(x) \to G(x) + f(x_0)$ במ''ש לכל. מצד שני, $f_n \to f$ שואף נקודתית (ובפרט במ''ש, ובכך סיימנו את החלק הראשון בהוכחה), וגבול נקודתי הוא יחיד, כלומר $f(x) = G(x) + f(x_0)$. נפעיל את המשפט היסודי של $G$ ונקבל $G'(x) = g(x)$ כלומר $f'(x) = g(x)$ כנדרש. 
	\end{proof}
	\begin{Theorem}[גרסה חזקה יותר של המשפט הקודם]
		יהיו $f_n, g \co [a, b] \to \R$. נניח: 
		\begin{itemize}
			\item $f_n \in C^{1}([a, b])$
			\item $f_n' \to g$ במ''ש
			\item קיימת $x_0 \in [a, b]$ כך ש־$(f_n(x_0))^{\infty}_{n = 1}$ מתכנסת ל־$c \in \R$. 
		\end{itemize}
		אז הפונקציה: 
		\[ f_n\overset{u}{\to} f(x) := c + \int^{x}_{x_0} \dequad \ g(t)\dt \]
		וכן $\frac{d}{\dx}f(x) = g(x)$ (נוסף על התכנסות במ''ש של $f_n$ ל־$f$). 
	\end{Theorem}
	בניגוד למשפט הקודם, עתה ''שחזרנו`` את הפונקציה $f$ בהינתן גבול נקודתי אחד, בלי שהיא הייתה נתונה לנו. 
	\begin{proof}
		נשאר כתרגיל בעבור הקורא
		
		\textbf{רמז: }יש הוכחה ברשימת של שירי
	\end{proof}
	\rmark{צריך לזכור שלהתכנסות במידה שווה מספיק (בעיקר ש־) שהנגזרות מתכנסות במ''ש. }
	\textbf{דוגמה: }הטור $\frac{1}{1 - x} = 1 + x + x^{2} + \cdots$. אם נגזור איבר־איבר נקבל $\frac{1}{(1 - x)^{2}} = 1 + 2x + 3x^{2} + 4x^{3} + \cdots$. ומכאן אפשר להראות התכנסות במ''ש של הפונקציה הראשונה (להשלים הוכחה). 
	
%	TODO
	''נורא שימושי``. 
	
	\subsection*{טורי פונקציות}
	\defi{יהיו $f, u_N \to I \to \R$ ($I$ קבוצה כללית). אז \textit{טור הפונקציות $\sumninf u_n(x) = f(x)$ מתכנס ל־$f$ במ''ש ב־$I$} אם סדרת הסכומים החלקיים $f_n(x)= \sumnko u_k(x)$ מתכנסת ל־$f$ במ''ש ב־$I$. }
	\noti{מקובל להוסיף את האות $u$ מעל השוויון (כך, $\overset{u}{=}$) כדי להדגיש ההתכנסות במ''ש. }
	\defi{תחת התנאים לעיל, נאמר ש־$\sumninf u_n(x) = f(x)$ \textit{נקודתית} אם לכל $\ag \in \R$, מתקיים $\sumninf u_n(\ag) = f(\ag)$ כטור מספרים. }
	\begin{Theorem}[קריטריון קושי להתכנסות טור פונקציות במידה שווה]
		מתקיים $\sumninf u_n(x)= f(x)$ במ''ש אמ''מ לכל $\eg > 0$ קיים $N_0 \in \N$ כך שלכל $m \ge n < N_0$ מתקיים: 
		\[ \forall x \in I \co \sof{\sum_{k = n}^{m}u_k(x)} < \eg \]
	\end{Theorem}
	\exe{ההוכחה טרוויאלית כי פשוט מעבירים את זה לסוכמים חלקיים ומפעילים את קריטריון קושי. ההוכחה כתרגיל. }
	\rmark{שימו לב, הכמתים בסדר שונה מקריטריון קושי מחדו''א 1א: קיים $N_0$ יחיד כך שלכל $x$ מתקיים (הא''ש לעיל), במקום שלכל $x$ קיים $N_0$ שמקיים (הא''ש לעיל). }
	\exe{אם $\sumninf u_n(x) = f(x)$ נקודתית, אז $u_n(x) \to 0$ נקודתית. אם $\sumninf u_n(x)$ מתכנס במ''ש, אז $u_n(x) \to 0$ במ''ש ב־$I$. }
	\rmark{זהירות! $\int^{\infty}_0 g(t) \dt$ מתכנס \textbf{לא גורר} ש־$g(t) \to 0$ (כאשר $t$ שואף לאינסוף). }
	\defi{נתונות $u_n \co I \to \R$. הטור $\sumninf u_n(x)$ מתכנס בהחלט במ''ש ב־$I$ אם $\sumninf \sof{u_n(x)}$ מתכנס במ''ש ב־$I$. }
	\theo{אם טור מתכנס בהחלט במ''ש, אז $\sumninf \sof{u_n(x)}$ מתכנס לכל $x$ וגם $\sumninf u_n(x)$ מתכנס במ''ש (הוכחה באמצעות קריטריון קושי במ''ש). \textbf{הכיוון השני לא עובד. }התכנסות בהחלט + התכנסות במ''ש לא גורר התכנסות בהחלט במ''ש. }
	\textbf{דוגמה. }נתבונן ב־$\log(1 + x) = \sum_{n = 1}^{\inft}(-1)^{n -1}\frac{x^{n}}{n}$ בקטע $[0, 1)$. היא מתכנסת במ''ש ומתכנסת בהחלט, אך לא בהחלט במ''ש. 
	
	\subsubsection*{משפטים מבניים וטרוויאלים}
	\theo{נניח $\sumninf u_n(x) =f(x)$ במ''ש ב־$I = [a, b]$. 
	\begin{enumerate}
		\item אם לכל $n$, $u_n$ רציפה אז $f(x)$ רציפה. 
		\item אם לכל $n$, $u_n$ אינטגרבילית אז $f(x)$ אינטגרבילית. במקרה זה מתקיים: 
		\[ \int^{b}_a \dequad \ f(x) \dx  = \sumninf \int^{b}_a \dequad \ u_n(x)\dx \]
		לפעמים יקרא ''אינטגרציה איבר־איבר``. 
	\end{enumerate}}
	\begin{Theorem}[דיני]
		אם $f, u_n$ רציפות וכן $\forall n \in \N \co u_n \ge 0$. אז אם $\sumninf u_n(x) = f(x)$ נקודתית אז $\sumninf u_n(x) = f(x)$ במ''ש. 
	\end{Theorem}
	
	
	
	
	
	
	
	
	\ndoc
\end{document}